% Regulatory Analysis subsection — standalone LaTeX fragment
% Intended to replace or expand \subsection{Regulatory Implications} in universal_impossibility.tex

\subsection{Regulatory Analysis}
\label{sec:regulatory-analysis}

We analyze three major regulatory frameworks through the lens of
Theorem~\ref{thm:universal}, showing that each implicitly demands
all three desiderata simultaneously---and that our impossibility result
clarifies how compliance should be structured.

\paragraph{EU AI Act Article~13 (Transparency).}
The EU AI Act \citep{euaiact2024} requires that high-risk AI systems
be designed with ``appropriate transparency'' and that users receive
``meaningful information about the logic involved'' in automated
decisions.  In the language of our framework, ``meaningful information''
implies that explanations must be \emph{faithful} (accurately reflecting
the model's actual reasoning) and \emph{decisive} (providing a definite
answer, not a hedge).  The qualifier ``appropriate'' is intended to hold
across model versions, deployment contexts, and regulatory audits---i.e.,
the explanation must also be \emph{stable}.
Theorem~\ref{thm:universal} shows that all three simultaneously is
impossible whenever the Rashomon property holds.  Concretely, two
equally-performant models of a credit scoring system may assign opposite
feature importance rankings; demanding a single ``meaningful'' ranking
that is faithful to both is a contradiction.

\emph{Resolution.}  Report $G$-invariant explanations (e.g., \DASH{}
attributions averaged over the Rashomon set) with explicit uncertainty
quantification.  Disclose \emph{which} property is sacrificed---typically
decisiveness---and why: the data genuinely underdetermines the
explanation.  Article~13 compliance is then achieved by providing
stable, faithful-in-expectation explanations with transparent
disclosure of ties.

\paragraph{NIST AI Risk Management Framework (MAP~1.5, MEASURE~2.6).}
The NIST AI RMF \citep{nistai2023} requires that ``interpretability and
explainability are defined, appropriate to the AI risk, and are
documented'' (MEASURE~2.6), and that risks from lack of interpretability
are mapped and managed (MAP~1.5).  The word ``appropriate'' signals a
context-dependent balancing of explanation quality---but the framework
provides no formal guidance on \emph{which} properties of an explanation
should take priority.  In practice, this implicitly assumes that
faithfulness, stability, and decisiveness are all simultaneously
achievable: one simply chooses the ``appropriate'' level of each.

\emph{Resolution.}  Our theorem shows that compliance should explicitly
specify \emph{which leg} of the trilemma is prioritized for each use
case.  For safety-critical applications (e.g., medical devices), one may
prioritize faithfulness and accept instability across model versions.
For auditable financial models, one may prioritize stability and accept
ties.  The NIST framework's ``define and document'' mandate is best
fulfilled by documenting the trilemma tradeoff explicitly in the system
card.

\paragraph{OCC SR~11-7 (Model Risk Management).}
SR~11-7 \citep{occ2011sr117} establishes three pillars for model
validation in banking: \emph{conceptual soundness} (the model's
theoretical basis must be justified), \emph{outcomes analysis} (model
performance must be monitored across versions and time), and
\emph{effective challenge} (independent reviewers must produce specific,
actionable findings).  These three requirements map directly onto our
trilemma: conceptual soundness demands \emph{faithfulness} (the
explanation reflects the model's actual structure), outcomes analysis
across model versions demands \emph{stability}, and effective challenge
demands \emph{decisiveness} (a reviewer cannot ``effectively challenge''
a model if the explanation is a set of ties).  SR~11-7 thus demands all
three for every banking model subject to validation.

\emph{Resolution.}  For model validation, use ensemble explanations
(\DASH{}) that are stable and partially faithful---faithful in
expectation over the Rashomon set---satisfying conceptual soundness and
outcomes analysis.  Document decisiveness limitations explicitly in
model documentation: where the data underdetermines the explanation,
state so and provide the $G$-invariant confidence interval rather than a
point estimate.  Effective challenge is then directed at the
\emph{modeling choices} (why this Rashomon set? why this loss threshold?)
rather than at explanation artifacts that are mathematically
underdetermined.

\paragraph{Synthesis.}
Our theorem does not imply that explainability requirements are
wrong---it implies they must be stated as \emph{tradeoffs}, not
absolutes.  A regulation that demands ``faithful, stable explanations''
is achievable (drop decisive).  A regulation that demands ``faithful,
decisive explanations'' is achievable (accept instability across
versions).  A regulation that demands all three is asking for the
mathematically impossible.  The $G$-invariant resolution provides the
best achievable compliance: stable, faithful-in-expectation, with
explicit ties where the data underdetermines the explanation.
Regulators who adopt this framing---requiring disclosure of \emph{which}
desideratum is relaxed and \emph{why}---will produce more honest, more
useful, and more legally defensible AI governance than those who demand
the impossible.
